\documentclass[12pt,a4paper,oneside]{report}
\usepackage[a4paper,left=30mm,right=20mm,top=25mm,bottom=25mm]{geometry}
\usepackage{fontspec}\usepackage{polyglossia}\setmainlanguage{mongolian}\setmainfont{Noto Sans}
\usepackage{graphicx,booktabs,longtable,array,hyperref,setspace,listings,xcolor}
\onehalfspacing\hypersetup{colorlinks=true,linkcolor=blue,urlcolor=blue}
\begin{document}
\begin{titlepage}\centering{\large МОНГОЛ УЛСЫН ИХ СУРГУУЛЬ\par}\vspace{.5cm}{\large ПРОГРАММ ХАНГАМЖИЙН ХӨТӨЛБӨР\par}\vspace{2.5cm}{\LARGE\bfseries LABSMART\par}\vspace{.5cm}{\Large Лабораторийн төхөөрөмжийн удирдлага, ашиглалтын бүртгэл ба ирээдүйн хэрэгцээний таамаглал бүхий веб систем\par}\vspace{1cm}Төслийн тайлан\vfill Гүйцэтгэсэн: С. Мөнх-Оргил (23B1NUM0124)\\Хамтран гүйцэтгэсэн: Н. Мөнхцэцэг\\Удирдагч багш: \rule{6cm}{.4pt}\vfill Улаанбаатар хот, 2026 он\end{titlepage}
\pagenumbering{roman}\chapter*{Хураангуй}LabSmart нь лабораторийн төхөөрөмжийн бүртгэл, хүсэлт, QR олголт/буцаалт, засвар, аудит, тайлан болон ирээдүйн эрэлтийн таамаглалыг нэгтгэсэн веб систем юм. ML pipeline нь session × төхөөрөмжийн төрөл түвшинд requested\_qty target ашиглаж, дүрмийн baseline, Linear Regression болон Random Forest загварыг хугацааны дарааллаар харьцуулна. Synthetic demo өгөгдөл нь бодит нарийвчлалын нотолгоо биш.\chapter*{Abstract}LabSmart centralizes laboratory equipment inventory, requests, QR checkout and return, maintenance, audit, reporting, and short-term demand forecasting. Its pipeline uses requested quantity as the target and compares a rule baseline with Linear Regression and Random Forest. Synthetic data is used only for software validation.\tableofcontents\listoffigures\listoftables\clearpage\pagenumbering{arabic}
\chapter{Удиртгал}\section{Асуудлын үндэслэл}Лабораторийн төхөөрөмжийн бүртгэл, олголт, буцаалт, засвар болон ирээдүйн хэрэгцээг нэг дор удирдах шаардлага нь LabSmart төслийн үндсэн үндэслэл юм.\section{Зорилго ба зорилт}Одоо байгаа веб системийн бүтэц, өгөгдөл, функциональ дээр тулгуурлан төхөөрөмжийн удирдлага болон эрэлтийн таамаглалыг сайжруулна.
\chapter{Шаардлагын шинжилгээ}Оюутан каталог, хүсэлт, QR, хугацааг ашиглана. Багш/ажилтан хүсэлт шийдэж, олголт, буцаалт, засварыг хянана. Админ хэрэглэгч, тохиргоо, аудит болон системийн модулиудыг удирдана.
\chapter{Архитектур ба өгөгдлийн сан}Одоогийн багц Python HTTP сервер, SQLite, HTML/CSS/JavaScript интерфэйс болон Python ML pipeline ашиглана. Гол хүснэгтүүд: users, equipment, loans, lab\_sessions, requests, demand\_history, model\_runs, predictions, recommendations, maintenance, audit\_events, notifications, app\_settings.
\chapter{Хэрэгжүүлэлт}Төхөөрөмжийн каталог, хайлт/шүүлтүүр, хүсэлт батлах, QR-аар олгох/буцаах, мэдэгдэл, засвар, аудит болон CSV тайланг нэг урсгалд холбосон.\chapter{Эрэлтийн таамаглал}Нэг мөр нь session × төхөөрөмжийн төрөл; target нь олголтоос өмнөх requested\_qty. Cutoff-оос өмнөх feature ашиглаж, rule baseline, Linear Regression болон Random Forest-ийг хугацааны дарааллаар үнэлнэ. MAE, RMSE, Bias хэмжүүр хадгална.\chapter{Аюулгүй байдал ба хэрэглэгчийн туршлага}Password hashing, HttpOnly/SameSite cookie, CSRF token, rate limit, server-side role checks болон input validation нь үндсэн хамгаалалт. Production-д HTTPS, persistent database, backup/restore болон monitoring нэмнэ.\chapter{Туршилт ба хязгаарлалт}README-д 21 integration test-ийн suite гэж тэмдэглэсэн боловч upgrade notes-д 20/21 PASS гэсэн зөрүү бий. Иймээс яг ашиглах хувилбар дээр тестийг дахин ажиллуулж баталгаажуулна. Synthetic history нь бодит ML нарийвчлалын нотолгоо биш.\chapter{Дүгнэлт}LabSmart нь лабораторийн төхөөрөмжийн бүртгэл, хүсэлт, QR олголт/буцаалт, засвар, аудит, тайлан болон эрэлтийн таамаглалыг нэгтгэнэ. Цаашид бодит өгөгдөл, PostgreSQL, scheduled ML batch, persistent deployment, backup/restore, monitoring болон хэрэглэгчийн туршилт шаардлагатай.\begin{thebibliography}{9}\bibitem{template}МУИС, \textit{Дипломын ажил бичих загвар}, Overleaf template PDF, 2025. \url{https://www.overleaf.com/latex/templates/mkut-diplom-dadlagyn-azhlyn-tailan-bichikh-zagvar/jwmtvxdgjbfs.pdf}\bibitem{labguide}\textit{Лабораторийн тоног төхөөрөмжийн хэрэглээний таамаглал ба ухаалаг зөвлөмж бүхий веб удирдлагын систем}, төслийн удирдамж.\bibitem{mlguide}\textit{Лабораторийн төхөөрөмжийн эрэлтийг машин сургалтаар таамаглах удирдамж}, ML удирдамж.\bibitem{pandas}pandas documentation, \url{https://pandas.pydata.org/docs/getting_started/install.html}\bibitem{sklearn}scikit-learn documentation, \url{https://scikit-learn.org/stable/modules/model_evaluation.html}\end{thebibliography}\appendix\chapter{Ажиллуулах заавар}\begin{lstlisting}[language=bash]python -m pip install -r requirements.txt\ncd app\npython server.py --demo\end{lstlisting}\end{document}
